\documentclass[11pt]{article}

\usepackage[a4paper,margin=1in]{geometry}
\usepackage{amsmath}
\usepackage{amssymb}
\usepackage{amsfonts}
\usepackage{graphicx}
\usepackage{booktabs}
\usepackage{bm}
\usepackage[hidelinks]{hyperref}

% --- notation helpers ---
\newcommand{\E}{\mathbb{E}}
\newcommand{\Prob}{\mathbb{P}}
\newcommand{\eps}{\varepsilon}
\newcommand{\dd}{\,d}
\DeclareMathOperator{\Tr}{Tr}
\DeclareMathOperator*{\supop}{sup}

\title{Pairs Trading and Statistical Arbitrage Strategies}
\author{}
\date{}

\begin{document}
\maketitle

\noindent\textbf{Note.} This is a self-contained write-up of the theory behind the pairs-trading and
statistical-arbitrage strategy implemented in this repository. The text is paraphrased and the eight
figures are \emph{generated from the repository's own code} (\texttt{make\_book\_figures.py}, which
calls the optimal-stopping solver \texttt{book\_opt.py} together with the backtest's models). As a
correctness check, the entry/exit trigger levels plotted in Figure~\ref{fig:11_6} agree to four
decimals with the trigger table given in Table~\ref{tab:triggers}.

\section{Pairs Trading and Statistical Arbitrage Strategies}

\subsection{Introduction}

How well a trading algorithm performs usually comes down to how accurately it can anticipate price
moves, and forecasting a basket of assets is typically easier than forecasting any one of them in
isolation. \textbf{Pairs trading} exploits exactly this: rather than betting on a single name, one
trades a portfolio built as a linear combination of two assets, chosen so that their \emph{joint}
behaviour is more predictable than either leg alone. The crucial ingredient is the way the two assets
move together. Consider, for instance, two assets whose spread (the gap between their prices) tends to
follow a recognisable pattern and only departs from it temporarily. A pairs-trading rule then makes
money by wagering that such departures eventually revert to their usual level. Strategies of this kind
are often grouped under the heading of \textbf{statistical arbitrage} (StatArb). They are not genuine
arbitrages---which earn more than the risk-free rate at zero risk---but instead lean on the
\emph{typical} dynamics of prices, so they do carry risk.

\begin{figure}[ht]
  \centering
  \includegraphics[width=\textwidth]{figures/fig_11_1.pdf}
  \caption{Two assets that move together over a trading window: (left) each midprice divided by its
  mean; (right) the resulting co-integration factor, shown with its mean (dashed) and $\pm 2$
  standard-deviation bands (dash-dotted). A real co-integrated crypto pair (LTC, UNI) is used here as
  the illustrative example.}
  \label{fig:11_1}
\end{figure}

Figure~\ref{fig:11_1} illustrates the idea with two assets that broadly track one another in the same
direction. Going long one leg and short the other produces a combined position that is calmer and
easier to forecast than either asset by itself. A sharper version of the idea works with the
co-integration factor of the two prices, a quantity of the form $\eps_t = A\,S^{(1)}_t + B\,S^{(2)}_t$
whose coefficients $A$ and $B$ are fitted to the data. When the mean-reverting pattern persists, a
portfolio holding $A$ units of the first asset and $|B|$ units short of the second should fluctuate
around the average level of $\eps_t$. The trading rule is then straightforward: buy the portfolio when
it looks cheap and unwind once it richens, or sell it when it looks expensive and unwind once it
cheapens.

The remainder is organised as follows. Section~\ref{sec:adhoc} covers simple rules that place fixed
bands around the equilibrium level; Section~\ref{sec:optimal} derives the bands that are optimal in a
precise sense; and Section~\ref{sec:sta} treats the case where it is the \emph{drifts} of a basket of
assets, rather than their prices, that are co-integrated.

\subsection{Ad Hoc Bands}\label{sec:adhoc}

The most basic way to harvest the factor's mean-reversion is to set bands one standard deviation either
side of the equilibrium level (taken to be zero here): buy one unit of the portfolio when the lower band
is touched and sell one unit when the upper band is touched. Having opened a position, the rule then
waits for the factor to come back close to equilibrium---say within $1/10$ of a standard deviation---and
closes out there.

\begin{figure}[ht]
  \centering
  \includegraphics[width=0.8\textwidth]{figures/fig_11_2.pdf}
  \caption{One simulated path of the co-integration factor together with the strategy's position
  (inventory) and its running book value, under a standard-deviation band rule.}
  \label{fig:11_2}
\end{figure}

The three panels of Figure~\ref{fig:11_2} show, respectively, a simulated factor path, the inventory as
the rule opens and closes trades, and the cumulative marked-to-market book value. Operationally, the
rule waits until the path crosses an outer band, opens the matching position betting on reversion, and
then carries the portfolio, whose value tracks the factor one-for-one. Its book value is
\begin{equation}
  BV_t = X_t + \beta_t\left( A\,S^{(1)}_t + B\,S^{(2)}_t \right),
\end{equation}
where $X_t$ is accumulated cash and $\beta_t \in \{-1,0,1\}$ is the number of portfolio units held. The
position is unwound once the factor returns to the inner band, banking a profit roughly equal to the
spacing between the outer and inner bands.

In the path shown the strategy happens to finish flat, but that is not guaranteed: a position opened
near the end of the horizon may never revert in time, leaving the trader holding inventory and exposed
to losses. Widening the bands makes leftover inventory more likely; it also raises the profit per
completed round trip while reducing how many outer-to-inner crossings occur over the horizon.

\begin{figure}[ht]
  \centering
  \includegraphics[width=\textwidth]{figures/fig_11_3.pdf}
  \caption{P\&L histograms over 10{,}000 simulated paths for the naive rule at several band widths:
  (a) $0.25\times$, (b) $0.5\times$, (c) $1.0\times$, (d) $2.0\times$ the standard deviation. The
  Sharpe ratio rises and then falls as the band grows, and the widest band produces a clearly
  multimodal distribution.}
  \label{fig:11_3}
\end{figure}

Figure~\ref{fig:11_3} collects the profit-and-loss (P\&L) over 10{,}000 simulated paths, each time
forming the factor and trading as above. The Sharpe ratio---mean P\&L over its standard
deviation---climbs as the band widens and then turns down again. The histograms are also multimodal:
because each completed round trip earns about one band's width, the P\&L piles up near whole multiples
of the band, and the weight on each multiple reflects how likely that many round trips are within the
horizon. The mass is not concentrated solely at one band width because, near the end of the horizon,
the factor may not have reverted, forcing the trader to close at less than a full band's profit---or
even at a loss.

\subsection{Optimal Band Selection}\label{sec:optimal}

We now look for the genuinely optimal entry and exit rule by framing it as an optimal-stopping problem
in which the trader executes a single round trip and faces no fixed end date. Take the portfolio value
(the co-integration factor) to follow Ornstein--Uhlenbeck dynamics,
\begin{equation}
  d\eps_t = \kappa(\theta - \eps_t)\dd t + \sigma \dd W_t,
\end{equation}
with $W_t$ a standard Brownian motion, $\kappa$ the reversion speed, $\theta$ the long-run level, and
$\sigma$ the volatility. The payoff criterion for closing a long position is
\begin{equation}
  H_+^{(\tau)}(t,\eps) = \E_{t,\eps}\!\left[ e^{-\rho(\tau - t)}(\eps_\tau - c) \right],
  \qquad
  H_+(t,\eps) = \supop_{\tau} H_+^{(\tau)}(t,\eps),
\end{equation}
in which $c$ is the cost of unwinding and $\rho>0$ is the discount rate (effectively an impatience
knob---larger $\rho$ drags the exit boundary toward the mean). The criterion for opening a long
position is
\begin{equation}
  G_+^{(\tau)}(t,\eps) = \E_{t,\eps}\!\left[ e^{-\rho(\tau - t)}\big(H_+(\tau,\eps_\tau) - \eps_\tau - c\big) \right],
  \qquad
  G_+(t,\eps) = \supop_{\tau} G_+^{(\tau)}(t,\eps);
\end{equation}
on entry the trader pays $\eps + c$ and in return holds the exit option, worth $H_+$.

Because the OU process is stationary, the value functions carry no explicit time dependence. The
dynamic programming principle then forces $H$ and $G$ to obey the coupled variational inequalities
\begin{align}
  \max\big\{(\mathcal{L} - \rho)H_+(\eps)\,;\,(\eps - c) - H_+(\eps)\big\} &= 0,\\
  \max\big\{(\mathcal{L} - \rho)G_+(\eps)\,;\,(H_+(\eps) - \eps - c) - G_+(\eps)\big\} &= 0,
\end{align}
where $\mathcal{L}$ is the generator of the factor process,
\begin{equation}
  \mathcal{L} = \kappa(\theta - \eps)\,\partial_\eps + \tfrac{1}{2}\sigma^2\,\partial_{\eps\eps}.
\end{equation}

\subsubsection{The Optimal Exit Problem}

The inequality governing $H$ is structurally like a perpetual call, and is solved by first finding the
fundamental solutions of
\begin{equation}\label{eq:111}
  (\mathcal{L} - \rho)F(\eps) = 0, \tag{11.1}
\end{equation}
which we label $F_\pm(\eps)$. We then write $H_+(\eps) = A\,F_+(\eps) + B\,F_-(\eps)$ on the
continuation set ($\eps < \eps^*$) and $H_+(\eps) = \eps - c$ on the exercise set ($\eps > \eps^*$),
imposing value matching $H_+(\eps^*) = \eps^* - c$ and smooth pasting $\partial_\eps H_+(\eps^*)=1$.
Explicitly, the fundamental solutions are
\begin{equation}
  F_+(\eps) = \int_0^\infty u^{\frac{\rho}{\kappa}-1}\,
    e^{-\sqrt{\frac{2\kappa}{\sigma^2}}(\theta - \eps)u - \frac12 u^2}\dd u,
  \qquad
  F_-(\eps) = \int_0^\infty u^{\frac{\rho}{\kappa}-1}\,
    e^{+\sqrt{\frac{2\kappa}{\sigma^2}}(\theta - \eps)u - \frac12 u^2}\dd u.
\end{equation}
A short calculation shows $F_+$ is positive, increasing and convex, whereas $F_-$ is positive,
decreasing and convex. Since $H$ has to decay to zero as $\eps \to -\infty$, the coefficient $B$
vanishes and
\begin{equation}
  H_+(\eps) = A\,F_+(\eps)\,\mathbf{1}_{\eps<\eps^*} + (\eps - c)\,\mathbf{1}_{\eps\ge\eps^*}.
\end{equation}
The two boundary conditions give $A\,F_+(\eps^*)=\eps^*-c$ and $A\,F_+'(\eps^*)=1$; dividing one by the
other pins down the exit level as the unique root of
\begin{equation}\label{eq:112}
  (\eps^* - c)\,F_+'(\eps^*) = F_+(\eps^*), \tag{11.2}
\end{equation}
with $A = (\eps^*-c)/F_+(\eps^*)$, so that
\begin{equation}
  H_+(\eps) = \frac{F_+(\eps)}{F_+(\eps^*)}(\eps^* - c)\,\mathbf{1}_{\eps<\eps^*}
    + (\eps - c)\,\mathbf{1}_{\eps\ge\eps^*}.
\end{equation}

\begin{figure}[ht]
  \centering
  \includegraphics[width=\textwidth]{figures/fig_11_4.pdf}
  \caption{Exit value function $H_+$ (solid where the trader waits, dashed where she exercises) and the
  exit triggers (circles), for varying reversion speed $\kappa$ and discount rate $\rho$
  ($\theta=0$, $\sigma=0.5$, $c=0.01$).}
  \label{fig:11_4}
\end{figure}

Figure~\ref{fig:11_4} traces $H_+$ as $\kappa$ and $\rho$ change. Faster reversion pulls the trigger
levels lower, because the factor is tugged back toward equilibrium more forcefully. A larger discount
rate has the same effect, since it makes the trader less willing to wait for future payoffs.

\subsubsection{The Optimal Entry Problem}

With the exit rule in hand, the entry problem is another perpetual American option, this time with
exercise value $H_+(\eps) - \eps - c$. Expecting $G_+$ to fall with $\eps$, we set
\begin{equation}
  G_+(\eps) = D\,F_-(\eps)\,\mathbf{1}_{\eps>\eps_*} + \big(H_+(\eps) - \eps - c\big)\,\mathbf{1}_{\eps\le\eps_*}.
\end{equation}
Matching value and slope gives $D\,F_-(\eps_*) = H_+(\eps_*) - \eps_* - c$ and
$D\,F_-'(\eps_*) = H_+'(\eps_*) - 1$, and the entry trigger is the root of
\begin{equation}\label{eq:113}
  \big(H_+(\eps_*) - \eps_* - c\big)F_-'(\eps_*) = \big(H_+'(\eps_*) - 1\big)F_-(\eps_*). \tag{11.3}
\end{equation}
As expected $\eps_* < \eps^*$, and
\begin{equation}
  G_+(\eps) = \frac{F_-(\eps)}{F_-(\eps_*)}\big(H_+(\eps_*) - \eps_* - c\big)\mathbf{1}_{\eps>\eps_*}
    + \big(H_+(\eps) - \eps - c\big)\mathbf{1}_{\eps\le\eps_*}.
\end{equation}

\begin{figure}[ht]
  \centering
  \includegraphics[width=\textwidth]{figures/fig_11_5.pdf}
  \caption{Entry value function $G_+$ (solid where the trader waits, dashed where she exercises) and the
  entry triggers (circles), for varying $\kappa$ and $\rho$ ($\theta=0$, $\sigma=0.5$, $c=0.01$).}
  \label{fig:11_5}
\end{figure}

Figure~\ref{fig:11_5} plots $G_+$ for the same parameter sweeps. Stronger reversion pushes the entry
trigger up, closer to the equilibrium level. Comparing Figures~\ref{fig:11_5} and~\ref{fig:11_4} shows
that entry and exit are not mirror images about the mean: discounting nudges the entry to sit nearer
the equilibrium than the exit does.

\subsubsection{Double-Sided Optimal Entry--Exit}

Allow now both directions---the trader may open either long or short and later close optimally. The two
exit criteria are
\begin{equation}
  H_+^{(\tau)}(t,\eps) = \E_{t,\eps}\big[e^{-\rho(\tau-t)}(\eps_\tau - c)\big],
  \qquad
  H_-^{(\tau)}(t,\eps) = \E_{t,\eps}\big[e^{-\rho(\tau-t)}(-\eps_\tau - c)\big],
\end{equation}
with $H_\pm = \supop_\tau H_\pm^{(\tau)}$. The long side $H_+$ is unchanged from before, so only $H_-$
needs attention; it decreases in $\eps$ and tends to zero as $\eps \to \infty$:
\begin{equation}
  H_-(\eps) = A\,F_-(\eps)\,\mathbf{1}_{\eps>\eps^*_-} - (\eps + c)\,\mathbf{1}_{\eps\le\eps^*_-}.
\end{equation}
Value matching and smooth pasting deliver $A\,F_-(\eps^*_-) = -(\eps^*_- + c)$ and
$A\,F_-'(\eps^*_-) = -1$, so the short-exit trigger solves
\begin{equation}\label{eq:114}
  (\eps^*_- + c)\,F_-'(\eps^*_-) = F_-(\eps^*_-). \tag{11.4}
\end{equation}
Collecting both sides, the exit value functions read
\begin{align}
  H_+(\eps) &= \frac{F_+(\eps)}{F_+(\eps^*_+)}(\eps^*_+ - c)\,\mathbf{1}_{\eps<\eps^*_+}
    + (\eps - c)\,\mathbf{1}_{\eps\ge\eps^*_+},\\
  H_-(\eps) &= -\frac{F_-(\eps)}{F_-(\eps^*_-)}(\eps^*_- + c)\,\mathbf{1}_{\eps>\eps^*_-}
    - (\eps + c)\,\mathbf{1}_{\eps\le\eps^*_-}.
\end{align}
On the entry side, the value function $G$ is a mix of both fundamental solutions,
$G(\eps) = A\,F_+(\eps) + B\,F_-(\eps)$, on the band between the two entry triggers, with $A$ and $B$
fixed by value matching and smooth pasting at each trigger. A clean symmetry emerges: each optimal
entry coincides with the optimal exit of the opposite side---entering long happens at the short-exit
level, and vice versa.

\begin{figure}[ht]
  \centering
  \includegraphics[width=\textwidth]{figures/fig_11_6.pdf}
  \caption{Two-sided problem: value functions $H_+$ (solid) and $H_-$ (dashed) with their optimal
  triggers, for several $\kappa$ ($\rho=0.01$, $\sigma=0.5$, $\theta=1$, $c=0.01$). Circles mark the
  long-exit levels and squares the short-exit levels; the accompanying table lists the computed values.}
  \label{fig:11_6}
\end{figure}

Figure~\ref{fig:11_6} reports the value functions and triggers at $\rho=0.01$, $\sigma=0.5$,
$\theta=1$, $c=0.01$. The exit levels are not symmetric about $\theta$, and each entry level matches the
opposite side's exit level. The numerical triggers are listed below.

\begin{table}[ht]
  \centering
  \begin{tabular}{ccccc}
    \toprule
    & \multicolumn{2}{c}{exit triggers} & \multicolumn{2}{c}{entry triggers}\\
    \cmidrule(lr){2-3}\cmidrule(lr){4-5}
    $\kappa$ & $\eps^*_+$ (long) & $\eps^*_-$ (short) & $\eps_{*+}$ (long) & $\eps_{*-}$ (short)\\
    \midrule
    0.5 & 1.9537 & $-0.4060$ & $-0.4060$ & 1.9537\\
    1.0 & 1.7460 & $-0.1815$ & $-0.1815$ & 1.7460\\
    2.0 & 1.5744 & $-0.0740$ & $-0.0740$ & 1.5744\\
    4.0 & 1.4367 & $-0.0410$ & $-0.0410$ & 1.4367\\
    \bottomrule
  \end{tabular}
  \caption{Optimal exit and entry triggers for the two-sided problem. The repository's optimal-stopping
  solver reproduces these numbers exactly (see \texttt{make\_book\_figures.py}).}
  \label{tab:triggers}
\end{table}

\section{Co-integrated Log Prices with Short-Term-Alpha}\label{sec:sta}

A second route is to co-integrate the \emph{drift terms} of a basket of assets instead of the prices
directly.

\subsection{Model Setup}

Let a basket of risky assets have price vector $Y = (Y^1_t,\dots,Y^n_t)_{0\le t\le T}$ evolving as
\begin{equation}\label{eq:115}
  \frac{dY^k_t}{Y^k_t} = \delta_k\,\alpha_t \dd t + \sum_{i=1}^n \sigma_{ki}\dd W^i_t,
  \qquad
  \alpha_t = a_0 + \sum_{i=1}^n a_i \log Y^i_t, \tag{11.5}
\end{equation}
where $W=(W^1,\dots,W^n)$ are independent Brownian motions. The instantaneous covariance of assets $i$
and $j$ is $\Omega_{ij} = \sum_{k=1}^n \sigma_{ik}\sigma_{kj}$, making $\bm\sigma$ the Cholesky factor
of $\bm\Omega = \bm\sigma\bm\sigma'$, which we take to be invertible. Here $\alpha_t$ plays the role of
the co-integration factor (a `short-term-alpha' term). Differentiating it gives
\begin{equation}\label{eq:116}
  d\alpha_t = \kappa(\theta - \alpha_t)\dd t + \bm a'\bm\sigma \dd \bm W_t, \tag{11.6}
\end{equation}
with
\begin{equation}
  \kappa = -\sum_{k=1}^n a_k \delta_k = -\bm\delta'\bm a,
  \qquad
  \theta = \frac{\sum_{k=1}^n a_k \Omega_{kk}}{2\sum_{k=1}^n a_k\delta_k}
    = \frac{1}{2}\frac{\Tr(\bm A\bm\Omega)}{\bm\delta'\bm a},
\end{equation}
where $\bm A = \mathrm{diag}(\bm a)$. Mean-reversion requires $\bm\delta'\bm a < 0$. In terms of
$\bm Z_t = \log \bm Y_t$ the log-prices follow a (singular) vector-autoregression,
$d\bm Z_t = (\bm c - \bm B\bm Z_t)\dd t + \bm\sigma\dd\bm W_t$.

Trading is assumed frictionless and the objective is utility of terminal wealth. Writing $\bm\pi =
(\pi^0_t,\dots,\pi^n_t)$ for the dollar amounts held in the riskless and risky assets, the share count
in asset $k$ is $m^k_t = \pi^k_t / Y^k_t$ and wealth is
$X^{\bm\pi}_t = \sum_{k=0}^n \pi^k_t = \sum_{k=0}^n m^k_t Y^k_t$, evolving as
\begin{equation}
  dX^{\bm\pi}_t = (\bm\pi_t'\bm\delta)\,\alpha_t \dd t + \bm\pi_t'\bm\sigma \dd\bm W_t.
\end{equation}

\subsection{The Agent's Optimisation Problem}

Take an exponential-utility agent, $u(x) = -e^{-\gamma x}$, with value function
\begin{equation*}
  H(t,x,\bm y) = \supop_{\bm\pi\in\mathcal A}\E_{t,x,\bm y}\big[-\exp(-\gamma X^{\bm\pi}_T)\big].
\end{equation*}
Its Hamilton--Jacobi--Bellman equation is
\begin{equation}
  \partial_t H + \alpha\bm\delta'\bm D_{\bm y}H + \tfrac12 \bm D^{\bm\Omega}_{\bm{yy}}H
    + \supop_{\bm\pi}\Big\{(\bm\pi'\bm\delta)\alpha\,\partial_x H
    + \tfrac12(\bm\pi'\bm\Omega\bm\pi)\partial_{xx}H + \bm\pi'\bm\Omega\bm D_{\bm{xy}}H\Big\} = 0,
\end{equation}
with terminal data $H(T,x,\bm y) = -e^{-\gamma x}$. Completing the square in $\bm\pi$ yields the
feedback-form optimal holdings,
\begin{equation}\label{eq:117}
  \bm\pi^* = -\frac{\bm\Omega^{-1}\bm\delta\,\alpha\,\partial_x H + \bm D_{\bm{xy}}H}{\partial_{xx}H}, \tag{11.7}
\end{equation}
and the nonlinear PDE
\begin{equation}\label{eq:118}
  \partial_t H + \alpha\bm\delta'\bm D_{\bm y}H + \tfrac12 \bm D^{\bm\Omega}_{\bm{yy}}H
    - \frac{\mathcal{L}'H\,\bm\Omega^{-1}\mathcal{L}H}{2\,\partial_{xx}H} = 0. \tag{11.8}
\end{equation}

\subsection{Solving the DPE}

Guided by the Merton problem, try the ansatz
\begin{equation}
  H(t,x,\bm y) = -\exp\Big\{-\gamma\big(x + h(t,\alpha)\big) + \sum_{i=1}^n a_i \log y^i\Big\},
\end{equation}
with $\alpha = a_0 + \sum_i a_i \log y^i$ and $h(T,\alpha)=0$. Inserting it, the nonlinear pieces cancel
and $h$ is left obeying a plain linear PDE,
\begin{equation}\label{eq:1110}
  \partial_t h - \tfrac12\Tr(\bm A\bm\Omega)\,\partial_\alpha h + \tfrac12(\bm a'\bm\Omega\bm a)\,
    \partial_{\alpha\alpha}h + \frac{\bm\delta'\bm\Omega\bm\delta}{2\gamma}\,\alpha^2 = 0, \tag{11.10}
\end{equation}
with $h(T,\alpha)=0$. The optimal holdings reduce to
\begin{equation}\label{eq:1111}
  \bm\pi^* = \frac1\gamma\big(\bm\Omega^{-1}\bm\delta\big)\alpha - \bm a\,\partial_\alpha h. \tag{11.11}
\end{equation}
The leading term reproduces the classical Merton allocation (with the stochastic drift
$\bm\delta\alpha_t$), while the second term, along $\bm a$, is the adjustment introduced by
co-integration.

\subsection{A Probabilistic Interpretation}

Under the risk-neutral measure $\Prob^*$ generated by the market price of risk
$\bm\lambda_t = -\bm\sigma^{-1}\bm\delta\,\alpha_t$, the factor loses its level drift and becomes
\begin{equation}\label{eq:1113}
  d\alpha_t = -\tfrac12\Tr(\bm A\bm\Omega)\dd t + \bm a'\bm\sigma\dd\bm W^*_t. \tag{11.13}
\end{equation}
A Feynman--Kac representation then expresses $h$ as
\begin{equation}\label{eq:1114}
  h(t,\alpha) = \E^*_{t,\alpha}\!\left[\frac{\bm\delta'\bm\Omega\bm\delta}{2\gamma}\int_t^T \alpha_s^2\dd s\right], \tag{11.14}
\end{equation}
so the extra value from trading this co-integrated basket comes from the expected future integral of
$\alpha^2$, growing with both volatility (via $\bm\Omega$) and the loading $\bm\delta$.

\subsection{Explicit Construction of the Optimal Investment Strategy}

Evaluating the expectation makes $h(t,\alpha)$ quadratic in $\alpha$, and the optimal holdings take the
closed form
\begin{equation}\label{eq:1116}
  \bm\pi^* = \frac1\gamma\left\{(\bm\Omega^{-1}\bm\delta)\,\alpha
    + (\bm\delta'\bm\Omega\bm\delta)\Big[\tfrac12\tau\,\alpha + \tfrac14\Tr(\bm A\bm\Omega)\tau^2\Big]\bm a\right\},
  \qquad \tau = T-t. \tag{11.16}
\end{equation}
The first piece is the Merton portfolio (the assets drift at $\bm\delta\alpha$); the second, directed
along $\bm a$, is the co-integration adjustment, and it shrinks as $\tau\to 0$---near the end of the
horizon the trader stops chasing the short-term-alpha signal.

\subsection{Numerical Experiments}

As an illustration we take a three-asset example with
\[
  \bm\delta = (1,1,0)', \quad \bm a = (-1,0,1)', \quad a_0 = 0,
\]
\[
  \bm\sigma = \begin{pmatrix} 0.2000 & 0 & 0\\ 0.0375 & 0.1452 & 0\\ 0.0250 & 0.0039 & 0.0967 \end{pmatrix}
  \;\Rightarrow\;
  \bm\Omega = \begin{pmatrix} 0.04 & 0.0075 & 0.005\\ 0.0075 & 0.0225 & 0.0015\\ 0.005 & 0.0015 & 0.01 \end{pmatrix},
\]
volatilities $\{2\%, 1.5\%, 1\%\}$, and starting prices $Y^1_0=11.10$, $Y^2_0=12.00$, $Y^3_0=11.00$.
With this choice of $\bm a$, only the first and third assets enter the factor, while $\bm\delta$ lets
the factor act on the drifts of the first two assets only.

\begin{figure}[ht]
  \centering
  \includegraphics[width=\textwidth]{figures/fig_11_7.pdf}
  \caption{A single simulated path: the co-integrated asset prices, the factor $\alpha_t$, the implied
  optimal holdings $\bm\pi^*_t$, and the wealth process $X_t$.}
  \label{fig:11_7}
\end{figure}

\begin{figure}[ht]
  \centering
  \includegraphics[width=0.75\textwidth]{figures/fig_11_8.pdf}
  \caption{Distribution of terminal P\&L for the optimal strategy across many simulated paths.}
  \label{fig:11_8}
\end{figure}

Figure~\ref{fig:11_7} displays one realisation of the prices, factor, optimal holdings, and wealth,
while Figure~\ref{fig:11_8} gives the histogram of terminal P\&L over many simulations.

\section{Bibliography and Selected Readings}

Elliott, Van Der Hoek \& Malcolm (2005); Tourin \& Yan (2013); Leung \& Li (2014); Leung \& Ludkovski
(2011).

\end{document}
